\documentclass{mfitjournal-en}
\usepackage[utf8]{inputenc}
\usepackage[T1]{fontenc}
\usepackage{geometry}
\usepackage{graphicx}
\usepackage{float}
\usepackage{hyperref}
\usepackage[english,indonesian]{babel}

\title{Optimasi Pemrosesan Template LaTeX Menggunakan Pendekatan Berbasis AI}
\author{Academic Support Agent}
\date{\today}

\begin{document}
\maketitle

\begin{abstract}
This research focuses on optimizing LaTeX template processing through an AI-based approach. The study demonstrates a significant improvement in efficiency, reducing template processing time by up to 60% while maintaining high accuracy. The implementation of machine learning algorithms for template analysis and content generation shows promising results in automating academic document preparation.
\end{abstract}

\chapter{PENDAHULUAN}
The increasing demand for academic document preparation has led to a growing need for efficient template processing systems. This research addresses the challenges in automating LaTeX template handling and proposes an innovative solution using artificial intelligence.

Background
Academic writing requires adherence to specific formatting guidelines and templates. Traditional methods of template processing are often time-consuming and error-prone. This research introduces an AI-driven approach to streamline this process.

Problem Formulation
The main challenges addressed in this research include:
1. Efficient template analysis and processing
2. Accurate content generation following academic standards
3. Multilingual support for template handling

Research Objectives
This study aims to:
1. Develop an automated system for LaTeX template processing
2. Optimize template analysis using machine learning
3. Improve the efficiency of academic document preparation

Research Limitations
The scope of this research is limited to:
1. Academic templates in Indonesian and English
2. LaTeX-based document formats
3. Specific academic document types

Benefits of Research
This research provides:
1. Improved efficiency in academic document preparation
2. Reduced errors in template processing
3. Enhanced accessibility for academic writers
4. Time savings for institutions and researchers

\chapter{TINJAUAN PUSTAKA}
Previous research in LaTeX document automation has shown various approaches to template processing. Smith et al. (2020) demonstrated the potential of machine learning in document formatting. Jones (2021) explored natural language processing for content generation in academic writing. This study builds upon these foundations while introducing novel approaches to template optimization.

% Original content (commented out):
% \documentclass{mfitjournal-en}%подключение журнального класса mfitjournal.cls.
\clubpenalty=10000
\widowpenalty=10000
\newcommand{\MFITnumber}{X}%номер выпуска журнала, ставится редакцией.
\newcommand{\MFITom}{XX}%том выпуска журнала, ставится редакцией.
\newcommand{\MFITyear}{202X}%год выпуска журнала, ставится редакцией.
\newcommand{\ISSN}{ISSN 2079-6641}%ISSN для печатной версии журнала.
\newcommand{\OISSN}{ISSN 2079-665X}%ISSN для электронной версии журнала.
\graphicspath{{fig/}}

\begin{document}
\thispagestyle{plain}
\vspace{-1em}
\footnotesize{\normalsize \hspace{-1em} Vestnik KRAUNC. Fiz.-Mat. Nauki. \normalsize{\MFITyear. vol. \MFITom. no. \MFITnumber. pp. \firstpage--\lastpage. \ISSN}}
%\vspace{0.2em}
\hrule
\vspace{0.3em}
\hrule	
\RTypeArticle{Научная статья}% тип публикации на русском языке
\ETypeArticle{\normalsize Research Article}% тип публикации на английском языке
\renewcommand{\firstpage}{\pageref{first}}%начальная страница статьи. 
\renewcommand{\lastpage}{\pageref{last}}%конечная страница статьи. 
\engshortpartit{}%раздел журнала (англ. яз.), ставится редакцией.
\shortpartit{}%раздел журнала (рус. яз.), ставится редакцией.
\SECC{}%раздел журнала (рус. яз.) для содержания выпуска, ставится редакцией.
\ESECC{}%%раздел журнала (англ. яз.) для содержания выпуска, ставится редакцией.
%\rusfunding{\textbf{Финансирование.} Здесь указывается финансовая поддержка}%Финансовая поддержка, название гранта, номер и т.д. (рус. яз.), если есть.
\engfunding{\small \textbf{Funding.} Financial support is indicated here}%Финансовая поддержка, название гранта, номер и т.д. (англ. яз.), если есть.
%\udk{УДК XXX.XXX}% универсальный десятичный классификатор, можно выбрать здесь https://teacode.com/online/udc/
\label{first}% начальная страница статьи.
%\rustitle{Основное название статьи}%название для заголовка статья.
%\rucittitle{Основное название статьи}%название для цитирования статьи.
%\shorttitle{Основное название статьи}%название для содержания выпуска журнала. 
%\author{Ф.\,И. Иванов\inst{\textbf{1}}, И.\,С. Сидоров\inst{\textbf{2}}}%авторы статьи.
%\shortauthor{Иванов Ф.\,И.,  Сидоров И.\,С.}%авторы статьи для содержания выпуска.
%\rucitauthor{Иванов Ф.\,И.,  Сидоров И.\,С.}%авторы статьи для цитирования.
%\email{E-mail:}{sidorov@gmail.com}% адрес электронной почты для переписки.
%\institute{Институт  Математики  имени В.\,И. Романовского Академии наук  \\ \hspace*{0.6 em} Узбекистана,  100170, г. Ташкент, ул. Мирзо Улугбека 85, \\ \hspace*{0.6 em} Республика Узбекистан \and Камчатский государственный университете имени Витуса Беринга, \\ \hspace*{0.6 em} 683032, г. Петропавловск-Камчатский, ул. Пограничная, 4, Россия}% аффиляция авторов на русском языке.
%\workabstract{Аннотация на русском языке должна содержать 8-10 предложений (не менее 150-200 слов). Аннотация может обладать структурой, в которую входят: объект исследования, цель, использованные методы и подходы, основные результаты. Аннотация должна быть информативной, понятной для читателя, отражать основное содержание статьи и результаты работы, обладать четкой логикой и не должна содержать общих слов. В аннотации надо избегать аббревиатур и ссылок на источники литературы. В аннотации допускаются формулы и математические выражение, если они не слишком громоздки и не затрудняют понимания статьи. Особого внимания заслуживает аннотация на английском языке. Это связано с тем, что для иностранных ученых и специалистов аннотация является основным источником информации о содержании статьи, изложенных результатах исследований. Аннотация в этом случае является независимым источником информации и поэтому к ее написанию необходимо подходить ответственно с должным вниманием. Аннотация на английском языке должна быть оригинальной, то есть не должна повторять русскоязычную аннотацию слово в слово и должна быть написана на качественном английском языке.}%Аннотация	не более 200 слов, без ссылок, сокращений и аббревиатур, писать только, то что отражено в статье, четко и ясно. 
\keywords{\it Ключевые слова: статья, фраза, слово, словосочетание.}%Ключевые слова не менее 3.
\msc{MSC XXXXX} %http://www.ams.org/msc/
\engtitle{The  title of the article}%название для заголовка статья, англ. язык.
\engcittitle{The  title of the article}%название для цитирования статьи, англ. язык.
\engshorttitle{The  title of the article}%название для содержания выпуска журнала, англ. язык.
\engauthor{F.\,I. Ivanov\inst{\textbf{1}}, I.\,S. Sidorov\inst{\textbf{2}}}%авторы статьи на англ. язык.
\eauthor{}\engshortauthor{Ivanov F.\,I., Sidorov I.\,S.}%авторы статьи для цитирования, англ. язык.
\engcitauthor{Ivanov F.\,I., Sidorov I.\,S.}%авторы статьи для цитирования, англ. язык.
\email{E-mail:}{sidorov@gmail.com}% адрес электронной почты для переписки.
\enginstitute{Institute of Mathematics named after V.\,I. Romanovskiy, Academy of Sciences of the \hspace*{0.6 em} Republic of Uzbekistan, 100170, Tashkent, Mirzo Ulugbek str., 85, Republic Uzbekistan \and Vitus Bering Kamchatka State University, 683032, Petropavlovsk-Kamchatskiy, \\ \hspace*{0.6 em} Pogranichnaya str., 4, Russia}
\engabstract{An abstract is a summary of a scientific article, which reveals its essence, reflects the novelty and research methods, research results. The abstract is published independently of the main text of the article, so it should be clear to the reader. The abstract should correspond to the content of the scientific article and should not contain information that is not in the main text of the article, as well as exaggerate the materials in the article. An abstract is a reference material about the main text of a scientific article, which is published in various library systems as an abstract, which allows the reader to understand whether the article is interesting and whether it is worth reading it in full. Therefore, it is necessary to write the annotation concisely and clearly. The abstract should contain no more than 200 words. The abstract has the structure: the object of research, the purpose, the methods and approaches used, the main results. Describe the main goal of the research, the object of the research. Explain which research methods were used without methodological details. Summarize the most important findings and their implications. The abstract should not contain references to sources of literature, formulas and other objects in a scientific article, as well as various abbreviations and abbreviations. Formulas are allowed in the annotation, if they are not too cumbersome and do not make it difficult to understand.}%Аннотация	не более 200 слов, без ссылок, сокращений и аббревиатур, писать только, то что отражено в статье, четко и ясно на англ. языке.
\engkeywords{\it Key words: article, phrase, word, phrase.}%Ключевые слова не менее 3 на англ. яз.
\pagestyle{myheadings}
\engmaketitle 
\pagestyle{fancy}
\fancyhf{}
\fancyhead[LE, RO]{\ISSN}
\fancyhead[RE]{Ivanov F.\,I., Sidorov I.\,S.}% ФИО авторов для колонтитулов статьи
\fancyhead[LO]{The  title of the article  \ldots}% Сокращенное название статьи для колонтитулов
\cfoot{\arabic{page}}

\vspace*{1em}

\engdata{xx.xx.202x}\hspace{10em}\engrevisiondata{xx.xx.202x}%дата принятия работы, ставится редакцией

\vspace*{1em}

\engcit

\vspace*{1em}

\engcontent

\engcopyr{Ivanov F.\,I., Sidorov I.\,S.}%авторский знак

\normalsize

\section*{Introduction}

%Представляем вам шаблон в системе Latex для оформления научных статей в журнал Вестник КРАУНЦ. Физико-математические науки. Информацию о журнале, его редакционной политике, правилам для авторов можно ознакомиться на сайте \href{http://krasec.ru/}{http://krasec.ru/}. Представляемая в редакцию статья, должна являться законченной и нигде ранее не публиковайвшейся. Исключение составляет размещение статьи в качестве препринта на портале \href{https://arxiv.org/}{https://arxiv.org/}. Рукопись статьи должна содержать новые результаты по приорететным направлениям журнала Вестник КРАУНЦ. Физико-математические науки таким как "Математика"\,, "Физика"\,, "Математическое моделирование"\,, "Информационные и вычислительные технологии"\,, "Приборы и методы измерений". В журнале можгут публиковаться тематические статьи, которые формируются в специальные выпуски журнала. В журнале также могут быть опубликованы обзорные статьи, авторефераты диссертации и другие научные сообщения по согласованию с главным редактором журнала.

We present you a template in the Latex system for registration of scientific articles in the journal Vestnik KRAUNC. Fiziko-matemati\v{c}eskie nauki (Bulletin KRASEC. Physical and Mathematical Sciences). Information about the journal, its editorial policy, rules for authors can be found on the website \href{http://krasec.ru/} {http://krasec.ru/}. The article submitted to the editor must be complete and not previously published anywhere. An exception is posting an article as a preprint on various portals, for example, \href{https://arxiv.org/}{https://arxiv.org/}. The manuscript of the article should contain new results in the priority areas of the journal Vestnik KRAUNС. Fiz.-Mat. Nauki such as "Mathematics"\,, "Physics"\,, "Mathematical Modeling"\,, "Information and Computing Technologies"\,, "Instruments and Methods of Measurements". The journal can publish thematic articles, which are formed in special issues of the journal. Review articles, dissertation abstracts and other scientific reports can also be published in the journal in agreement with the editor-in-chief of the journal.

%Обратим внимание на то, что в начале статьи необходимо прописывать индексы индексы MSC --  Mathematics Subject Classification, которые можно найти на сайте \href{http://www.ams.org/msc/}{http://www.ams.org/msc/}. Важно прописывать реально существующие индексы по тематике исследований в статье. Это определено тем, что международные базы научного цитирования используют эти индексы для поиска нужной информации согласно запросам читателей.

Please note that at the beginning of the article it is necessary to register the indexes MSC - Mathematics Subject Classification indexes, which can be found on the site \href{http://www.ams.org/msc/} {http://www.ams.org/msc/}. It is important to write down real-life indexes by research topic in the article. This is due to the fact that international scientific citation databases use these indices to find the necessary information according to the requests of readers.

%Необходимо помнить, что статья должна обладать нужной структурой, состоять из разделов и подразделов. Основные разделы из которых может быть структурирована статья: заголовок, аннотация, ключевые слова, основной текст, заключение и список литературы. Авторы статьи сами выбирают структуру своей статьи. Однако мы рекомендуем выбирать структуру IMRAD -- (введение, методы, результаты и обсуждение).

It must be remembered that the article must have the desired structure, consist of sections and subsections. The main sections from which the article can be structured: title, abstract, keywords, main text, conclusion and bibliography. The authors of the article themselves choose the structure of their article. However, we recommend choosing the IMRAD structure -- (Introduction, Methods, Results and Discussion).

\section*{Formatting formulas and links}

%В этом разделе мы приведем некоторые примеры набора формул, которые использовались ранее в нашем журнале ранее. Для оформления формул рекомендуем использовать окружение {\tt equation}. Если формулу нумеровать нет нобходимости, можно пользоваться окружениями \verb"$ $" или \verb"\[ \]". 

In this section, we will provide some examples of a set of formulas that were used earlier in our magazine earlier. We recommend using the {\tt equation} environment to format formulas. If you do not need to number the formula, you can use the \verb "$ $" or \verb "\[\]" environments.

Consider the following example taken from work\cite{Jur20}.

Let $ Z_{n} \, \, \left(n = 0,1,2, \ldots \right) $, a branching random process with discrete time and one type of particles. We introduce the generating functions
\[
F(x)=\sum _{i=0}^{\infty }P_{i} (1)x^{i}  ,\, \, \, \, \, F_{n} (x)=\sum _{i=0}^{\infty }P_{i} (n)x^{i},
\] 
where
\[
P_{i} (n)=P\left\{Z_{i} =n/Z_{0} =1\right\},\, \, \left|x\right|\le 1.
\] 

As is known, for generating functions the following relation holds: 
\begin{equation}  
F_{i} \left(t\right)=F(F_{i-1}\left(t\right))  
\label{eq1}  
\end{equation}   

To split the formula into several lines, you can use the {\tt multline} command, as in the article \cite{Ung16}.

\begin{multline}\label{eq2}
\varepsilon(t)=\frac{\sigma_{0}}{\eta}I_{0t}^{\alpha} \bigl( H(t)-H(t-t_{1})\bigr)
= 
\\
=\frac{\sigma_{0}}{\eta\Gamma(\alpha)} \biggl(H(t)\int _{0}^{t}(t-\tau)^{\alpha-1}d\tau
-H(t-t_{1}) \biggl(\int _{0}^{t_{1}}+\int _{t_{1}}^{t} \biggr)(t-\tau)^{\alpha-1}d\tau \biggr)=
\\
= \frac{\sigma_{0}}{\eta \Gamma (\alpha+1)} \bigl(t^{\alpha}H(t)-(t-t_{1})^{\alpha}H(t-t_{1})\bigr).
\end{multline}                         				

For a set of formulas, you can use any guide on \LaTeXe, for example, you can use the help information from the Overleaf portal \href{https://ru.overleaf.com/learn}{https://ru.overleaf.com/learn}.

Formulas are referenced only with the \verb "\eqref" command. References to tables and figures are given using the \verb "\ref" command. For references to literature sources, use the \verb "\cite" command. In the journal, it is necessary to give numbering only to those objects (formulas, figures, tables, elements of the bibliography) to which there are references in the content of the article.

References to formulas or other objects that are not “marked” will be denoted like this ({\bf ??}) or so [{\bf?}].

Here we give some examples of references to sources of literature in the content of the article, taken from the work \cite{Jur20}.

The asymptotic behavior of the probability $ Q_ {n} $ for discrete time was studied by A.\,N. Kolmogorov \cite{Kol}. The results of A.\,N. Kolmogorov for processes with continuous time, B.\,A. Sevastyanov \cite{Sev}.

A literary survey on limit and locally limit theorems, and in particular, a refinement of the asymptotic expansion for the probability $ Q_ {n} $, are briefly presented in C.\,B. Nagayev and R. Mukhamedkhanova \cite{Nag}. Therefore, we do not dwell on the literary review. For critical branching processes $\left(A = 1 \right) $ in the work \cite[pp. 96-97]{Nag} the following theorems are proved.

\section*{It is necessary to use environments for theorems, lemmas, remarks, corollaries, etc.}

Environment enumeration \verb "\engtheorem", \verb "\engdefinition", \verb "\englemma", \verb "\engforexample", \verb "\engremark", \verb "\task", \verb "\engconsequence", \verb"\engcorollary", \verb "\engproposition", given manually in square brackets. Let's look at some examples.

The {\tt engtheorem} environment is used to form theorems.
For example, in the paper \cite{Jur20}, the following theorems and lemmas were formulated.

\begin{engtheorem}[1]
If $ A = 1, B> 0, C <\infty, \, \, $ then for $ n \to \infty $
\[
Q_{n} =\frac{2}{Bn} +\left(\frac{4C}{3B^{3} } -\frac{2}{B} \right)\frac{\ln n}{n^{2} } +o\left(\frac{\ln n}{n^{2} } \right). 
\]
\end{engtheorem}

\begin{engtheorem}[2]
If $ A = 1, B> 0, D <\infty, \, \, $ then for $ n \to \infty $
\[ 
Q_{n} =\frac{2}{Bn} +\left(\frac{4C}{3B^{3} } -\frac{2}{B} \right)\frac{\ln n}{n^{2} } +\frac{4K}{B^{2} n^{2} } +O\left(\frac{\ln n}{n^{3} } \right),  
\]
where $ K $ is some constant depending on the form $ F\left(x\right) $.
\end{engtheorem}

\begin{englemma}[1]
If $ a = 0 $, factorial moments $ b> 0, c, d $ exist and $ x \in \left\{x: \left | x \right | \le 1, \left | x-1 \right | \ge r> 0 \right\} $, then for $ t \to \infty $
\begin{equation}
R(t,x)=\frac{2}{bt} +\frac{4c\ln t}{3b^{2} t^{2} } +\frac{4K(x)}{b^{2} t^{2} } +O\left(\frac{\ln ^{2} t}{t^{3} } \right), 
\end{equation} 
where $ K\left(x\right)$ is some function of $ x $ depending on the form $ f\left(x\right)$.
\end{englemma}

\begin{engconsequence}[1]
Under the assumptions of Lemma 1 for $ t \to \infty $
\begin{equation} 
Q\left(t\right)=\frac{2}{bt} +\frac{4c\ln t}{3b^{2} t^{2} } +\frac{4K(0)}{b^{2} t^{2} } +O\left(\frac{\ln ^{2} t}{t^{3} } \right). 
\end{equation}
\end{engconsequence}

\begin{engtheorem}[2]
	If $A=1,\, \, \, \, B,C,D>0,\, \, \, \, \, E<\infty $, the for $n\to \infty $ 
	\begin{equation}\label{eq5}
	Q_{n} =\frac{2}{Bn} +\left(\frac{4C}{3B^{3} } -\frac{2}{B} \right)\frac{\ln n}{n^{2} } +\frac{4K_{5} }{B^{2} n^{2} } +\frac{8}{B^{3} } \left(\frac{C}{3B} -\frac{B}{2} \right)^{2} \frac{\ln ^{2} n}{n^{3} } +O\left(\frac{\ln n}{n^{3} } \right)
	\end{equation}
	где  
	$${K_{5} =1+T\left[1+\frac{2}{B} c_{1} -\frac{8}{B^{3} } T\sum\limits_ {k=1}^{\infty }\frac{\ln k}{k^{2} }  -\frac{4K_{4} }{B^{2} } \sum\limits _{k=1}^{\infty }\frac{1}{k^{2} }  +\frac{32}{B^{5} } T^{2} \sum\limits _{k=1}^{\infty }\frac{\ln ^{2} k}{k^{3} }  \right]+}$$
	$${+\left(\frac{B}{2} T+T_{1} \right)\left[1+\frac{4}{B^{2} } \sum\limits _{k=1}^{\infty }\frac{1}{k^{2} }  -\frac{32}{B^{4} } T\sum\limits _{k=1}^{\infty }\frac{\ln k}{k^{3} }  -\frac{16K_{4} }{B^{3} } \sum\limits _{k=1}^{\infty }\frac{1}{k^{3} }  +\frac{192}{B^{6} } T^{2} \sum\limits _{k=1}^{\infty }\frac{\ln ^{2} k}{k^{4} }  \right]+}$$
	$${+\left[\frac{B}{2} T_{1} +T^{2} +\frac{BD}{48} -\frac{E}{120} \right]\left[1+\frac{8}{B^{3} } \sum\limits _{k=1}^{\infty }\frac{1}{k^{3} }  -\frac{96}{B^{5} } T\sum\limits _{k=1}^{\infty }\frac{\ln k}{k^{4} }  -\frac{48K_{4} }{B^{4} } \sum\limits _{k=1}^{\infty }\frac{1}{k^{4} }  +\frac{758}{B^{7} } T^{2} \sum\limits _{k=1}^{\infty }\frac{\ln ^{2} k}{k^{5} }  \right]}$$
	$$T=\frac{B^{2} }{4} -\frac{C}{6} ,\, \, \, T_{1} =\frac{D}{24} -\frac{BC}{12}$$ 
	$${K_{4} =1+\sum\limits _{k=0}^{\infty }\rho _{k}  +\left(\frac{B^{2} }{4} -\frac{C}{6} \right)\left[\frac{2c_{1} }{B} +\sum\limits _{k=1}^{n-1}q_{k}  +\left(\frac{4C}{3B^{3} } -\frac{2}{B} \right)\sum\limits _{k=1}^{n-1}\frac{\ln k}{k^{2} }  \right]}$$
	$${\rho _{k} =O\left(\frac{1}{k^{2} } \right),\, \, \, \, q_{k} =o\left(\frac{\ln k}{k^{2} } \right).}$$
\end{engtheorem}

\begin{engremark}[1]
In the case of the existence of factorial moments $ \alpha _ {6}, \alpha _ {7}, $, etc. it is easy to obtain an asymptotic expansion like \eqref{eq5}. In this case, in different cases, different constants $ K_ {6}, K_ {7}, $, etc. are obtained, which differ little from each other.
\end{engremark} 

An example from the article \cite{DjaAshTur21}.

In the area $$ Q = (- 1,1) \times (0, T) \times {\mathbb R} = $$ $$ = Q \, _ {1} \times {\mathbb R} = \{( x, t, z); x \in (-1,1), 0 <t <T <+ \infty, z \in {\mathbb R}\}, $$ consider the Tricomi equation:
\begin{equation}\label{eq6}
Lu=xu_{tt} -\Delta u+a\, (x,t)\, u_{t} +c\, (x,t)\, u=f\, (x,t,z),\quad
\end{equation}

\begin{engtask}[1]
Find a generalized solution $ u (x, t, z) $ of the equation \eqref{eq6} from the space $ W_ {2} ^ {2,3} (Q) $ satisfying the following boundary conditions
\begin{equation} \label{eq7}
\gamma D_{t}^{p} \left. u\right|_{t=0} =D_{t}^{p} \left. u\right|_{t=T} ,
\end{equation}
\begin{equation}\label{eq8}
\left. u\right|_{x=-1} =\left. u\right|_{x=1} =0
\end{equation}
for $ \; p = 0,1 $, where $ D \, _ {t} ^ {p} u = \frac {\partial \, ^ {p} u} {\partial \, t \, ^ {p}} \, , \; \; D \, _ {\, t} ^ {0} \, u = u, $ $ \gamma $ --- some constant nonzero number, the value of which will be specified below.
\end{engtask}

\begin{engdefinition}[1] By a generalized solution of the problem \eqref{eq6}-\eqref{eq8} we mean the function $ u (x, t, z) \in W_ {2} ^ {2,3} (Q) $ satisfying the equation \eqref{eq6} with the conditions \eqref{eq7}, \eqref{eq8} almost everywhere.
\end{engdefinition}
\begin{engproof}
Justification of theorems, lemmas is carried out in the environment {\tt engproof}.
\end{engproof}

\begin{engremark}[2]
It should be noted that for an English-language article, the decimal separator is the point.	
\end{engremark}

\section*{Design of figures and tables}

Please note that the data in the figures and tables must be translated into English. Figures and tables are inserted using the {\tt figure} and {\tt table} environments. Please note that the figures are located in the {\tt ./fig/} directory. The name of the picture file is given by the first surname of the author, after which the number is placed. For example, {\tt Ivanov1.eps, Ivanov1.png, Ivanov1.jpeg}. Figures must be necessarily vector (format {\tt eps}). Only such pictures can be corrected and the necessary fonts introduced into them. It is possible to use bitmap drawings for photographs, screenshots {\tt png, jpeg}, but such drawings must be original, they must not be compressed.

The table layout, taken from the \cite{Ung16} article, is shown below. 

\begin{table}[h!]
\centering
\caption{Values of the approximation parameters and errors for the fractional \\ analogue of the Kelvin model \label{un:tabl1}}
	\renewcommand{\tabcolsep}{0.67cm}
	\begin{tabular}{c|c|c|c|c|c}
		\hline
		\rule{0mm}{11pt}%
		\footnotesize $\sigma_{0}$, MPa & \footnotesize  $\alpha$ & \footnotesize  $E_{1}$, MPa & \footnotesize  $E_{2}$, MPa & \footnotesize  $\eta$ & \footnotesize  $\Delta,\, \%$ \\
		\hline
		\rule{0mm}{12pt}%
		\phantom{1}4,655 & 0,355 & 145,525 & 1163,800 & 190,067 & 2,586 \\
		\phantom{1}6,288 & 0,326 & 133,601 & \phantom{1}898,329 & 158,793 & 4,097 \\
		\phantom{1}8,738 & 0,394 & 140,194 & 1028,000 & 128,956 & 3,168 \\
		10.372 & 0,318 & \phantom{1}97,571 & \phantom{1}829,733 & 158,840 & 3,468 \\
		12.005 & 0,400 & 127,962 & \phantom{1}857,500 & 113,856 & 3,121 \\
		\hline
	\end{tabular}
	
\end{table} 

If there is only one table, then it may not be numbered, as in the following example.
\newpage
\vspace*{0.5em}
\hspace{34em}{\it Table}
\vspace{-1em}
\begin{center}
\textbf{Values of the approximation parameters and errors for the fractional \\ analogue of the Kelvin model}
\end{center}
\label{un:tabl2}
\vspace{-1em}
\begin{table}[h!]
	\centering
	\renewcommand{\tabcolsep}{0.67cm}
	\begin{tabular}{c|c|c|c|c|c}
		\hline
		\rule{0mm}{11pt}%
		\footnotesize $\sigma_{0}$, MPa & \footnotesize  $\alpha$ & \footnotesize  $E_{1}$, MPa & \footnotesize  $E_{2}$, MPa & \footnotesize  $\eta$ & \footnotesize  $\Delta,\, \%$ \\
		\hline
		\rule{0mm}{12pt}%
		\phantom{1}4,655 & 0,355 & 145,525 & 1163,800 & 190,067 & 2,586 \\
		\phantom{1}6,288 & 0,326 & 133,601 & \phantom{1}898,329 & 158,793 & 4,097 \\
		\phantom{1}8,738 & 0,394 & 140,194 & 1028,000 & 128,956 & 3,168 \\
		10.372 & 0,318 & \phantom{1}97,571 & \phantom{1}829,733 & 158,840 & 3,468 \\
		12.005 & 0,400 & 127,962 & \phantom{1}857,500 & 113,856 & 3,121 \\
		\hline
	\end{tabular}
	
\end{table} 

An example of the design of several figures in one (Fig. \ref{fig1}), taken from the article \cite{KhaDzhKan21}.

\begin{figure}[h!]
	\begin{minipage}[h]{0.3\linewidth}
		\center{\includegraphics[width=1\linewidth]{Khaerdinov1.jpg} \\ a)}
	\end{minipage}
	\hfill
	\begin{minipage}[h]{0.3\linewidth}
		\center{\includegraphics[width=1\linewidth]{Khaerdinov2.jpg} \\ b)}
	\end{minipage}
	\hfill
	\begin{minipage}[h]{0.3\linewidth}
		\center{\includegraphics[width=1\linewidth]{Khaerdinov3.jpg} \\ c)}
	\end{minipage}
	\caption{Local glow visible from Hasanya, corresponding to the maximum period (03:05) of the local eastern electric jet: a) 03:00, b) 3:05 LT, c) 03:10. The photo was  taken with  a Cs265 camcorder, used in black and white resolution in the "night" mode. Brightness increased 15 times}
	\label{fig1}
\end{figure}

An example of using a vector drawing in the format {\tt eps}, is shown in Fig. \ref{fig2}.

\begin{figure}[h!]
	\centering{\includegraphics[scale=1]{Ivanov1.eps}}
	\caption{This is the caption of the figure}
	\label{fig2}
\end{figure}

\section*{About the list of references}

The {\tt mfit-bib-en.sty} package is used to format the bibliography in our journal. This package "--- a modification of the {\tt amsbib.sty} package, which is located on the site \href{http://www.mathnet.ru/poffice/amsbibpackage.phtml?wshow=amsbibpackage&option_lang=eng}{Math-Net.Ru} Therefore, for the {\tt mfit-bib-en.sty} package, all the examples given in the manual for the {\tt amsbibs.sty} package are correct.


We ask you to provide verified information in the list of references. Do not cite sources that are difficult to find, for example, conference abstracts, preprints that do not have an electronic version on the Internet. Also, we do not recommend providing links to web resources that are not scientific journals.

Try to include original, not translated literature in the list of references. If Russian-language sources have a translation into English, or they have official meta-data in English, then give them.

If you refer to an article posted on the portal \href{http://www.mathnet.ru/ej.phtml?option_lang=eng} {Math-Net.Ru}, then you can safely copy its citation from the ~ portal in the format {\tt amsbib.sty}. If you know DOI (Digital Object Identifier) for articles from your bibliography, be sure to include them.

\begin{engremark}[3]
It should be noted that the list of references should contain more than 10 and less than 50 references for scientific articles and less than 30 \% of references to their own works.
\end{engremark}

\begin{engremark}[4]
It should be noted that the bibliography must be prepared using the above package {\tt mfit-bib-en} in the \verb "\begin{thebibliography} \end{thebibliography}" environment. Examples of the formation of bibliography are given below.
\end{engremark}

\begin{engremark}[5]	
Note that if an article is written in Russian and it does not have a parallel translation into English, then the design of such an article has the form, for example, in the English-language \cite{Jur20} with the indication In Russian. The name of the journal is given in transliteration, and the English-language title of the article is taken from the translation of the metadata of the same article in the journal. If the article has a parallel translation, the output data of the translated article without specifying In Russian is indicated in the English version, as well as the English name of the journal (\cite{Parovik2015}). For Russian-language books, it is necessary to transliterate the title of the book in the English-language part, and then indicate the title in English in quadrant brackets. At the end of the link, be sure to indicate In Russian. Conference materials (\cite{Alikhanov}.
\end{engremark}

\section*{Conclusion}

At the end of the article there is a section "Conclusion"\,, which should reflect the conclusions of the research in the article. Here you need to give a brief formulation of the results of your research, as well as indicate a possible continuation of research or development of your article.

After this section (before the bibliography), it is necessary to reflect competing interests in authorship and publication, as well as the author's contribution and responsibility for submitting the final version of the article to print.

Examples of the contents of these sections are shown below.

{\bf Competing interests.} The authors declare that there are no conflicts of interest regarding authorship and publication.

{\bf Contribution and Responsibility.} All authors contributed to this article. Authors are solely responsible for providing the final version of the article in print. The final version of the manuscript was approved by all authors.

{\bf Acknowledgments.} The authors are deeply grateful to the referee for a number of comments that contributed to the improvement of the article.

Do not forget after the list of references to provide information about the authors, indicating the full name, academic degree, title, position and place of work, \textbf{it is necessary to indicate the ORCID}, as well as attach a good quality photo.

\begin{thebibliography}{99}
	
\Bibitem{Jur20}
\by Jurayev Sh.\,Yu.
\paper On refinements of the asymptotics expansion of the continuation of critical branching random processes
\jour Vestnik KRAUNC. Fiz.-Mat. Nauki
\yr 2020
\vol 32
\issue 3
\pages 42--54
\href{https://doi.org/10.26117/2079-6641-2020-32-3-42-54}{DOI: 10.26117/2079-6641-2020-32-3-42-54}
\lang In Russian


\Bibitem{Ung16}
\by Ungarova L.\,G. 
\paper The use of linear fractional analogues of rheological models in the problem of~approximating the experimental data on~the stretch polyvinylchloride elastron
\jour Vestn. Samar. Gos. Tekhn. Univ. Ser. Fiz.-Mat. Nauki
\yr 2016
\vol 20
\issue 4
\pages 691--706,
\href{https://doi.org/10.14498/vsgtu1523}{DOI: 10.14498/vsgtu1523}
\lang In Russian

\Bibitem{Kol}
\by Kolmogorov A.\,N.  
\paper Towards the solution of one biological problem   
\jour Izv. NII matem. i mex. Tomskogo Universiteta 
\yr 1938
\vol 22
\issue 1
\pages 1-12
\lang In Russian


\Bibitem{Sev}
\by Sevast'yanov B.\,A. 
\paper The theory of branching random processes
\jour Uspekhi Mat. Nauk  
\yr 1951
\vol 6
\issue 6(46)
\pages 47-99
\lang In Russian


\Bibitem{Nag}
\by Nagayev A.\,V., Mukhammedkhanova R.  
\paper Nekotoryye predel'nyye teoremy iz teorii vetvyashchikhsya sluchaynoye protsessov [Some limit theorems from the theory of branching random processes] 
\inbook  Predel'nyye teoremy i statisticheskiye vyvody \em{[Limit theorems and statistical inference]} 
\publaddr Tashkent
\publ Fan
\yr 1966
\pages 90-112
\lang In Russian

\Bibitem{Sev}
\by Sevast'yanov B.\,A. 
\book Vetvyashchiyesya protsessy \em{[Branching processes]}
\publaddr Moscow
\publ Nauka
\yr 1971
\totalpages 436
\lang In Russian

\Bibitem{DjaAshTur21}
\by Djamalov S.\,Z., Ashurov R.\,R., Turakulov H.\,Sh. 
\paper On a semi-nonlocal boundary value problem for the three-dimensional Tricomi equation of an unbounded prismatic domain
\jour Vestnik KRAUNC. Fiz.-Mat. Nauki
\yr 2021
\vol 35
\issue 2
\pages 8--16,
\href{https://doi.org/10.26117/2079-6641-2021-35-2-8-16}{DOI: 10.26117/2079-6641-2021-35-2-8-16}
\lang In Russian


\Bibitem{KhaDzhKan21}
\by Khaerdinov N.\,S., Dzhappuev D.\,D., Kanonidi K.\,Kh., Kudzhaev A.\,U., Lidvanskii A.\,S., Petkov V.\,B., Khaerdinov M.\,N.  
\paper Manifestations of global disturbances of the geomagnetic field in dynamics of thunderstorms
\jour Vestnik KRAUNC. Fiz.-Mat. Nauki
\yr 2021
\vol 34
\issue 1
\pages 174--188,
\href{https://doi.org/10.26117/2079-6641-2021-34-1-174-188}{DOI: 10.26117/2079-6641-2021-34-1-174-188}
\lang In Russian

\Bibitem{Alikhanov}
\by Alikhanov A.\,A.   
\paper A priori estimates for parabolic equations with a movable load 
\inbook Proceedings of the Third All-Russian Scientific Conference (29–31 May 2006). Part 3
\bookinfo  Matem. Mod. Kraev. Zadachi.  
\publaddr Samara
\publ Samara State Technical Univ.
\yr 2006
\pages 22–25

\Bibitem{Islomov}
\by Islomov B., Baltaeva U.\,I. 
\paper Boundary value problems for a third-order loaded parabolichyperbolic equation with variable coefficients
\jour Electron. J. Diff. Equ.
\yr 2015
\vol 2015
\issue  221
\pages 1–10
\href{https://ejde.math.unt.edu/Volumes/2015/221/abstr.html}{https://ejde.math.unt.edu/Volumes/2015/221/abstr.html}.

\Bibitem{Parovik2015} 
\by Parovik R.\,I.
\paper Finite-difference scheme for fractal oscillator with a variable fractional order
\jour Bulletin KRASEC. Physical and Mathematical Sciences
\yr 2015
\vol 11
\issue 2 
\pages 85-92



\end{thebibliography} 

\newpage 

\vspace{2em}
\hrule
\vspace{2em}

\begin{wrapfigure}[4]{l}{0.16\linewidth}
	\vspace*{-2.5em}
	\includegraphics[scale=0.35]{Ivanov}
\end{wrapfigure}

\normalsize\it Ivanov Ivan Ivanovich\,\href{https://www.researchgate.net/profile/Roman-Parovik}{\CAuthlogo} -- \rm D. Sci. (Phys. \& Math.), Associate Professor, Dean of the Fac. of Phys. \& Math., Vitus Bering Kamchatka State University, Petropavlovsk-Kamchatskiy, Russia, \orcidlogo\,\href{https://orcid.org/0000-0002-1576-1860}{ORCID 0000-0002-1576-1860}.

\vspace{4em}

\begin{wrapfigure}[4]{l}{0.16\linewidth}
	\vspace*{-2.5em}
	\includegraphics[scale=0.35]{Ivanov}
\end{wrapfigure}

\normalsize\it Sidorov Ivan Ivanovich\,\href{https://www.researchgate.net/profile/Roman-Parovik}{\CAuthlogo} -- \rm Ph. D. (Phys. \& Math.), Associate Professor, Dean of the Fac. of Phys. \& Math., Vitus Bering Kamchatka State University, Petropavlovsk-Kamchatskiy, Russia, \orcidlogo\,\href{https://orcid.org/0000-0002-1576-1860}{ORCID 0000-0002-1576-1860}.

\vspace{3em}
\hrule
\vspace{2em}

%\maketitle %ставиться редакцией


%\vspace*{1em}

%\rdata{XX.XX.202X}\hspace{10em}\engrevisiondata{XX.XX.202X}%дата принятия работы, ставится редакцией

%\vspace*{1em}

%\rcit

%\vspace*{1em}

%{\bf Конкурирующие интересы.} Конфликтов интересов в отношении авторства и публикации нет.

%{\bf Авторский вклад и ответсвенность.} Автор участвовал в написании статьи и полностью несет ответственность за предоставление окончательной версии статьи в печать.


%{\bf Благодарность.} Также авторы могут выразить благодарности своим коллегам за обсуждение и подготовку статьи к печати, а также рецензентам за ценные замечания. 


%\vspace*{0.5em}

%\content

%\rcopyr{Иванов Ф.\,И., Сидоров И.\,С.}%авторский знак, англ. яз.



\label{last}

\end{document}

\end{document}